\section{Square grids}\label{app:square}

This section proves the lower bound in question~(4) of Section~\ref{sec:open}. Throughout, $n\ge2$, the grid is $n\times n$, $\mathcal R$ is the set of its rectangles, tilted ones included, and $L=1+\ln n$ with the natural logarithm. For a point $v$ we write $d_v$ for the number of rectangles that have $v$ as a corner. By Lemma~\ref{lem:threecorners}, two different rectangles share at most two corners, so the symmetric difference of their corner sets has four, six, or eight points.

\subsection{Pairs with two common corners}

\begin{lemma}[two common corners]\label{lem:compress}
Let $R_1\ne R_2$ be rectangles with exactly two common corners $p$ and $q$. If $\{p,q\}$ is a side of both rectangles, or a diagonal of both, then $C(R_1)\triangle C(R_2)$ is the corner set of a rectangle.
\end{lemma}

\begin{proof}
Suppose $\{p,q\}$ is a side of both. The other corners are $p+u,q+u$ for $R_1$ and $p+v,q+v$ for $R_2$, where $u\ne v$ are nonzero vectors perpendicular to $q-p$. The four points $p+u$, $q+u$, $p+v$, $q+v$ are distinct, since $p+u=q+v$ would make $v-u=p-q$ a nonzero vector perpendicular to itself. None of them is $p$ or $q$, since the common corners are exactly $p$ and $q$. They form a parallelogram with sides $q-p$ and $v-u$, which are nonzero and perpendicular, so they are the corners of a rectangle.

Suppose $\{p,q\}$ is a diagonal of both. Both rectangles have center $o=(p+q)/2$, and their other corners are $o\pm u$ for $R_1$ and $o\pm v$ for $R_2$ with $|u|=|v|=|p-q|/2$ and $v\ne\pm u$. The segments from $o-u$ to $o+u$ and from $o-v$ to $o+v$ have the same midpoint and the same length. So $o\pm u$ and $o\pm v$ are the four distinct corners of a rectangle.
\end{proof}

\subsection{Lattice points on circles}

Let $r_2(t)$ be the number of pairs $(a,b)$ of integers with $a^2+b^2=t$, and let $\tau(t)$ be the number of positive divisors of $t$.

\begin{lemma}[divisor bound]\label{lem:divisor}
For every $t\ge1$, $r_2(t)\le4\tau(t)$. For every $\varepsilon>0$ there is a constant $K_\varepsilon$ such that $\tau(t)\le K_\varepsilon t^\varepsilon$ for all $t\ge1$.
\end{lemma}

\begin{proof}
Both parts are classical. Hardy and Wright \cite[Sections~16.9 and~16.10]{hardywright2008} prove that $r_2(t)=4\bigl(d_1(t)-d_3(t)\bigr)$, where $d_i(t)$ is the number of divisors of $t$ congruent to $i$ modulo $4$, so $r_2(t)\le4d_1(t)\le4\tau(t)$. The bound $\tau(t)=O(t^\varepsilon)$ for every fixed $\varepsilon>0$ is proved in their Section~18.1.
\end{proof}

\subsection{Counting rectangles and obstructions}

Let $\mathcal A$ be the family of four-point sets consisting of all corner sets $C(R)$ with $R\in\mathcal R$, together with the mixed differences. A \emph{mixed difference} is a set $C(R_1)\triangle C(R_2)$ for a pair $R_1\ne R_2$ with exactly two common corners that form a side of one rectangle and a diagonal of the other, the case that Lemma~\ref{lem:compress} leaves out. Let $\mathcal B$ be the family of all six-point sets $C(R_1)\triangle C(R_2)$ for pairs $R_1\ne R_2$ with exactly one common corner.

\begin{lemma}[rectangle and obstruction counts]\label{lem:counts}
Let $n\ge2$ and $L=1+\ln n$.
\begin{enumerate}[(a)]
\item For every point $v$, $d_v\le8n^2L$, and $|\mathcal R|\le2n^4L$.
\item For every $\varepsilon>0$, $|\mathcal A|\le A_\varepsilon n^{4+2\varepsilon}L$ with $A_\varepsilon=2(1+16\cdot2^\varepsilon K_\varepsilon)$, and $|\mathcal B|\le32n^6L^2$.
\item If a blackout $S$ of the $n\times n$ grid contains no member of $\mathcal A$ and no member of $\mathcal B$, then $S$ is valid.
\end{enumerate}
\end{lemma}

\begin{proof}
(a) Let $R$ be a rectangle with corner $v$. Its two sides from $v$ are perpendicular nonzero integer vectors. Write the first as $k(a,b)$ with $k\ge1$ and $\gcd(a,b)=1$. The second is perpendicular to $(a,b)$, so it is $(-bl,al)$ for a real $l$, and $l$ is a nonzero integer because the vector is integral and $\gcd(a,b)=1$. Put $d=\max(|a|,|b|)\ge1$. There are $8d$ integer vectors $(a,b)$ with $\max(|a|,|b|)=d$, since they lie on the boundary of the square $[-d,d]^2$. Both sides fit in the grid, so $kd\le n-1$ and $|l|d\le n-1$, which leaves at most $(n-1)/d$ values of $k$ and $2(n-1)/d$ values of $l$. The ordered pair of sides determines $k$, $(a,b)$ and $l$, and $R$ gives two ordered pairs, so
\[
2d_v\le\sum_{d=1}^{n-1}8d\cdot\frac{n-1}{d}\cdot\frac{2(n-1)}{d}\le16n^2\sum_{d=1}^{n-1}\frac1d\le16n^2L.
\]
Each rectangle has four corners, so $4|\mathcal R|=\sum_vd_v\le n^2\cdot8n^2L$.

(b) Fix $R_1\in\mathcal R$ and one of its four sides $\{p,q\}$. A rectangle $R_2$ that has $\{p,q\}$ as a diagonal has center $(p+q)/2$ and a second diagonal of the same length, joining two grid points $p',q'$. The vector $p'-q'$ is an integer vector with $|p'-q'|^2=|p-q|^2=t$, where $t\le2(n-1)^2<2n^2$, and it determines $p'$ and $q'$. So there are at most $r_2(t)\le4K_\varepsilon(2n^2)^\varepsilon$ such $R_2$ by Lemma~\ref{lem:divisor}. Every mixed difference arises from some $R_1$, side, and $R_2$ in this way, so
\[
|\mathcal A|\le|\mathcal R|\bigl(1+4\cdot4K_\varepsilon2^\varepsilon n^{2\varepsilon}\bigr)\le2n^4L\bigl(1+16\cdot2^\varepsilon K_\varepsilon\bigr)n^{2\varepsilon}=A_\varepsilon n^{4+2\varepsilon}L .
\]
A member of $\mathcal B$ comes from an unordered pair of rectangles with a common corner $v$, so
\[
|\mathcal B|\le\sum_v\binom{d_v}2\le\frac{n^2}2\bigl(8n^2L\bigr)^2=32n^6L^2 .
\]

(c) Suppose $S$ is not valid, so some $R_1\ne R_2$ have the same presentation, and $C(R_1)\triangle C(R_2)\subseteq S$ by Lemma~\ref{lem:diff}. If $R_1$ and $R_2$ have two common corners, the symmetric difference is the corner set of a rectangle by Lemma~\ref{lem:compress} or a mixed difference, and in both cases it is a member of $\mathcal A$. If they have one common corner, it is a member of $\mathcal B$. If they have none, then $C(R_1)\subseteq S$ is a member of $\mathcal A$. In each case $S$ contains a member of $\mathcal A\cup\mathcal B$.
\end{proof}

Avoiding $\mathcal A$ forbids the one fully black rectangle that the game allows, which is harmless for a lower bound.

\subsection{The lower bound}

\begin{theorem}[random blackouts with deletions]\label{thm:alteration}
We have
\[
\liminf_{n\to\infty}\frac{f(n,n)\,(1+\ln n)^{2/5}}{n^{6/5}}\ge\frac5{6\cdot192^{1/5}}=0.2911\ldots,
\]
so for every $c<5/(6\cdot192^{1/5})$ we have $f(n,n)\ge c\,n^{6/5}(1+\ln n)^{-2/5}$ for
all large $n$. In particular
$f(n,n)-2n\to\infty$.
\end{theorem}

\begin{proof}
Put $a=192^{-1/5}$, $T=n^{6/5}L^{-2/5}$ and $\theta=a\,n^{-4/5}L^{-2/5}$, which lies in
$(0,1)$ since $a<1$, and select each point of the grid independently with probability
$\theta$. Let $X$ be the number of selected points. Let $Y$ be the number of members of
$\mathcal A$ all of whose four points are selected, and $Z$ the number of members of
$\mathcal B$ all of whose six points are selected. By linearity of expectation,
\begin{equation}\label{eq:expect}
\mathbb E(X-Y-Z)=n^2\theta-|\mathcal A|\theta^4-|\mathcal B|\theta^6 .
\end{equation}
Here $n^2\theta=aT$. With $\varepsilon=1/10$ and $A_0=A_{1/10}$, Lemma~\ref{lem:counts}(b)
gives $|\mathcal B|\theta^6\le32n^6L^2a^6n^{-24/5}L^{-12/5}=32a^6T$ and
$|\mathcal A|\theta^4\le A_0n^{21/5}La^4n^{-16/5}L^{-8/5}\le A_0a^4n^{-1/5}T$, as $L\ge1$. Since
$32a^5=\frac16$, \eqref{eq:expect} gives
$\mathbb E(X-Y-Z)\ge\bigl(\frac56a-A_0a^4n^{-1/5}\bigr)T$, so some selection $S_0$ has
$X-Y-Z$ at least this value.

Starting from $S_0$, delete one point from each member of $\mathcal A\cup\mathcal B$
still contained in the current set. At most $Y+Z$ deletions leave a set $S$ with
$|S|\ge X-Y-Z$ and no such member. Deletions cannot create an obstruction, since the
families are fixed. Lemma~\ref{lem:counts}(c) makes $S$ valid, so
$f(n,n)/T\ge\frac56a-A_0a^4n^{-1/5}$, tending to $5/(6\cdot192^{1/5})$. Finally
$T/n\to\infty$ gives the last claim.
\end{proof}

The value $a=192^{-1/5}$ maximizes $a-32a^6$, and every $0<\varepsilon<1/5$ would do in
place of $1/10$. This is the alteration method of Alon and
Spencer~\cite[Section~3.2]{alonspencer2016}.
